\documentclass[12pt]{article}
\usepackage[utf8]{inputenc}
\usepackage[T1]{fontenc}
\usepackage[margin=1in]{geometry}
\usepackage{setspace}
\usepackage{xcolor}
\usepackage[colorlinks=true, urlcolor=blue!60!black, linkcolor=blue!60!black]{hyperref}
\setlength{\parindent}{2em}
\setlength{\parskip}{1.25em}
\renewcommand{\baselinestretch}{1.0}
\begin{document}
\begin{center}
{\LARGE \textit{Saenz v. Roe}, 526 U.S. 489 (1999)\par}
\vspace{0.8em}
{\large Case Brief --- Four-Page Edition\par}
\end{center}
\vspace{1em}

\section{Memory Jogger}

Reviving the Fourteenth Amendment's Privileges or Immunities Clause, the Supreme Court held 7--2 that California violated newly arrived residents' constitutional right to travel by capping their first-year welfare benefits at the level of their prior states of residence, and that Congress could not authorize states to do so (\href{https://www.law.cornell.edu/supremecourt/text/98-97}{526 U.S. 489}).

\section{Detailed Case Facts}

California ranked sixth among the states in nominal welfare benefits: a mother and child received \$456 a month, against \$275 in Arizona, and a family of four \$673, against \$144 in Mississippi. To save an estimated \$10.9 million a year, the legislature in 1992 enacted Cal.\ Welf.\ \& Inst.\ Code \S~11450.03, limiting families who had lived in California for less than twelve months to the maximum grant of their prior state. The cap applied whether or not the family had received welfare before and whatever its motive for moving.

In 1996, Congress's welfare reform law (PRWORA) replaced AFDC with TANF block grants and, in 42 U.S.C. \S~604(c), expressly authorized states to apply another state's benefit rules to families resident for less than twelve months. California announced enforcement beginning April 1, 1997. That day Brenda Roe, from Oklahoma, and Anna Doe, from the District of Columbia, sued for a class of newcomers.

\section{Procedural History}

An earlier challenge, \textit{Green v. Anderson}, produced a preliminary injunction under \textit{Shapiro v. Thompson} and \textit{Zobel v. Williams} but was dismissed as unripe in 1995. After PRWORA, Judge Levi of the Eastern District of California certified the class and again enjoined enforcement, holding that PRWORA did not change the analysis. The Ninth Circuit affirmed, and the Supreme Court affirmed on May 17, 1999.

\section{Judicial Votes}

Affirmed, \textbf{7--2}.
\begin{itemize}
\item \textbf{Majority:} Justice Stevens, joined by Justices O'Connor, Scalia, Kennedy, Souter, Ginsburg, and Breyer.
\item \textbf{Dissents:} Chief Justice Rehnquist, joined by Justice Thomas; Justice Thomas, joined by Chief Justice Rehnquist.
\end{itemize}

\section{Holding}

The right to travel has three components: entering and leaving another state; being treated as a ``welcome visitor rather than an unfriendly alien'' while temporarily present, under Article IV, \S~2; and, for those who become permanent residents, being treated like other citizens of the new state, under the Fourteenth Amendment's Citizenship and Privileges or Immunities Clauses. Discrimination against bona fide new citizens by length or place of prior residence faces a standard ``surely no less strict'' than strict scrutiny, cannot rest on saving money or deterring poor migrants, and cannot be authorized by Congress. Remarks on ``readily portable'' benefits such as divorce and tuition are dicta.

\section{Analysis of Opinions}

\textbf{Majority (Stevens).} The cap erected no barrier to entry, and Article IV's protection of visitors against ``the disabilities of alienage'' does not govern settled residents. The third component rests on Section 1 of the Fourteenth Amendment: both Miller's majority and Bradley's dissent in the \textit{Slaughter-House Cases} agreed that a citizen may become a citizen of any state by residence, with the same rights as its other citizens. Because the classification itself penalizes that right, neither rational basis nor intermediate scrutiny applies. Deterring poor migrants is impermissible; the \$10.9 million in savings could be achieved by an across-the-board cut of 72 cents a month and cannot justify 45 subclasses of citizens, because the Citizenship Clause does not allow ``degrees of citizenship based on length of residence.'' Cash aid is consumed in-state, unlike portable benefits, so sham residency is unlikely. Neither Article I nor Section 5 lets Congress authorize Fourteenth Amendment violations.

\textbf{Rehnquist (dissenting).} The majority revived a clause used only once before, in \textit{Colgate v. Harvey}, which was overruled. Once settled, a person is ``no longer `traveling' in any sense of the word.'' The cap was a good-faith residency test like those upheld for tuition and divorce, and he called the portability distinction one ``border[ing] on the metaphysical.''

\textbf{Thomas (dissenting).} Tracing the phrase ``privileges and immunities'' from the 1606 Charter of Virginia through \textit{Corfield v. Coryell}, Thomas argued the Clause covers fundamental rights, not public benefits like welfare. Open to reconsidering \textit{Slaughter-House}, he warned that reviving the Clause without its original meaning risks ``yet another convenient tool for inventing new rights.''

\section{Examples: Future Applications}

\begin{enumerate}
\item \textbf{Same-Side (Unconstitutional): Housing Stipend Cap.} A state limits newcomers of eighteen months to the housing stipend of their prior state. Like \S~11450.03, it creates subclasses of citizens for a non-portable benefit.
\item \textbf{Same-Side (Unconstitutional): Loyalty Tax Exemption.} A homestead exemption grows by \$1,000 per decade of residence to reward ``long-term civic loyalty,'' creating degrees of citizenship barred by \textit{Saenz} and \textit{Zobel}.
\item \textbf{Opposite-Side (Constitutional): Elk License Fees.} Charging visitors \$400 and residents \$40 for elk licenses falls under Article IV, where recreational hunting is not fundamental (\textit{Baldwin}).
\item \textbf{Opposite-Side (Constitutional): Divorce Residency.} A one-year residency rule for divorce protects against sham residency for a ``readily portable'' decree (\textit{Sosna}).
\item \textbf{Fence-Sitter: Two-Year Tuition Rule.} A university denies in-state tuition for two years even to working, tax-paying adults who have clearly settled. \textit{Saenz} called tuition portable, but a rigid wait for bona fide residents may penalize new citizens.
\end{enumerate}

\section{Critique}

\textbf{Originalist and textualist views.} Some originalists praised the Court for breaking its long silence on the Clause after \textit{Slaughter-House}. Many echo Thomas: Stevens ignored Reconstruction history, and welfare entitlements sit uneasily with the \textit{Corfield} tradition. Kurt Lash and John Harrison debate whether the Clause protects national rights or equality among state citizens.

\textbf{Progressive and doctrinal views.} Progressives welcomed protection of poor migrants and the limit on congressional authorization. Laurence Tribe and Linda Bosniak warned that moving rights from Equal Protection, which covers ``any person,'' to a clause protecting citizens could weaken protection for noncitizens. The revival proved a one-off: in \textit{McDonald} only Thomas relied on the Clause.

\section{Key Quotations}

\begin{quote}
``The `right to travel' discussed in our cases embraces at least three different components. It protects the right of a citizen of one State to enter and to leave another State, the right to be treated as a welcome visitor rather than an unfriendly alien when temporarily present in the second State, and, for those travelers who elect to become permanent residents, the right to be treated like other citizens of that State.''
\end{quote}
\noindent --- Justice Stevens, Opinion of the Court, 526 U.S. at 500.

\begin{quote}
``Congress has no affirmative power to authorize the States to violate the Fourteenth Amendment and is implicitly prohibited from passing legislation that purports to validate any such violation.''
\end{quote}
\noindent --- Justice Stevens, Opinion of the Court, 526 U.S. at 508.

\begin{quote}
``Citizens of the United States, whether rich or poor, have the right to choose to be citizens `of the State wherein they reside.' U. S. Const., Amdt. 14, \S~1. The States, however, do not have any right to select their citizens.''
\end{quote}
\noindent --- Justice Stevens, Opinion of the Court, 526 U.S. at 510--511.

\end{document}
